\documentclass[12pt]{article}
\usepackage{amsfonts,amsthm,amsmath,amssymb,mathtools}
\usepackage{mathrsfs}
\usepackage{enumitem}
\usepackage{tikz-cd}
\usepackage{geometry}
\usepackage{mlmodern}
\usepackage{xcolor}
\usepackage{pgfplots}

\geometry{letterpaper, portrait, margin=1in}
\pgfplotsset{compat=1.18}

\setcounter{section}{0}

\renewcommand{\thesection}{\S\arabic{section}}
\renewcommand{\thesubsection}{\arabic{section}}
\newcommand{\dd}{\mathop{}\!\mathrm{d}}

\newtheorem{thm}{Theorem}
\newenvironment{theorem}[1]{
	\renewcommand{\thethm}{#1}
	\begin{thm}
		}{\end{thm}}

\newtheorem{lma}{Lemma}
\newenvironment{lemma}[1]{
	\renewcommand{\thelma}{#1}
	\begin{lma}
		}{\end{lma}}

\theoremstyle{definition}
\newtheorem{ex}{}
\newenvironment{exercise}[1]{
	\renewcommand{\theex}{#1}
	\begin{ex}
		}{\end{ex}}

\title{Homework Assignment 7}
\author{Connor Mooneyhan}
\date{October 17, 2025}

\begin{document}

\maketitle

I waited too long to start thinking about these. I'll be back on track next week.

\begin{ex}
	(2E.3)
	Give an example of a sequence of continuous functions $f_1, f_2, \dots$ from $[0, 1]$ to $\mathbb{R}$ that converges pointwise to a function $f : [0, 1] \to \mathbb{R}$ that is not a bounded function.
\end{ex}

\begin{ex}
	(2E.6)
	Suppose $(X, S, \mu)$ is a measure space with $\mu(X) < \infty$.
	Suppose $f_1, f_2, \dots$ is a sequence of $S$-measurable functions from $X$ to $\mathbb{R}$ such that $\lim_{k \to \infty}f_k(x) = \infty$ for each $x \in X$.
	Prove that for every $\varepsilon > 0$, there exists a set $E \in S$ such that $\mu(X \setminus E) < \varepsilon$ and $f_1, f_2, \dots$ converges uniformly to $\infty$ on $E$ (meaning that for every $t > 0$, there exists $n \in \mathbb{Z}^+$ such that $f_k(x) > t$ for all integers $k \geq n$ and all $x \in E$.)
\end{ex}

\begin{ex}
	(3A.3)
	Suppose $(X, S, \mu)$ is a measure space and $f : X \to [0, \infty]$ is an $S$-measurable function.
	Prove that \begin{equation*}
		\int f \dd \mu > 0 \text{ if and only if } \mu(\lbrace x \in X : f(x) > 0 \rbrace) > 0.
	\end{equation*}
\end{ex}

\begin{ex}
	(3A.4)
	Give an example of a Borel measurable function $f : [0, 1] \to (0, \infty)$ such that $L(f, [0, 1]) = 0$.
\end{ex}

\begin{ex}
	(3A.5)
	Verify the assertion that integration with respect to counting measure is summation.
\end{ex}
% \begin{proof}
% 	Suppose $\mu$ is counting measure on $\mathbb{Z}^+$ and $b_1, b_2, \dots$ is a sequence of nonnegative numbers.
% 	Let $b$ be the function from $\mathbb{Z}^+$ to $[0, \infty)$ defined by $b(k) = b_k$.
% 	Then \begin{align*}
% 		\int b \dd \mu & = \sup \left\lbrace \mathcal{L}(b, P) : \text{$P$ is an $S$-partition of $\mathbb{Z}^+$} \right\rbrace                                    \\
% 		               & = \sup \left\lbrace \sum_{j=1}^m \mu(A_j) \inf_{A_j} b : \text{$P = A_1, \dots, A_m$ is an $S$-partition of $\mathbb{Z}^+$} \right\rbrace \\
% 	\end{align*}
% \end{proof}

\end{document}
